\chapter{Conclusioni e sviluppi futuri}
\label{ch:conclusions}

\section{Sintesi rispetto alle questioni di ricerca}
\label{sec:summary-research-questions}

La \Cref{sec:objectives-research-questions} aveva formulato due questioni: quale vantaggio predittivo la letteratura attribuisca alla struttura relazionale tra criptovalute, e come quella struttura evolva nel passaggio da un mercato stabile a uno in crisi. I capitoli precedenti hanno costruito gli strumenti per affrontarle e ne hanno raccolto le risposte disponibili; questa sezione le mette insieme, senza aggiungere nulla di nuovo.

\subsection*{Prima questione: il vantaggio della struttura}

La risposta si articola su due piani, e conviene tenerli distinti perché hanno solidità diversa.

Sul piano teorico l'esito è netto. Un modello che tratta ogni asset come un sistema chiuso scarta informazione disponibile, e la Graph Convolutional Network è un modo ben fondato di recuperarla. Il \Cref{ch:from-cnn-to-gnn} ha mostrato che la sua regola di propagazione non è una formula ad hoc ma il punto d'arrivo di una catena di semplificazioni ciascuna motivata: filtri spettrali liberi, vincolo polinomiale — che regala la localizzazione dimostrata nella \cref{prop:laplacian-powers-locality} — base di Chebyshev, riduzione a un solo passo, rinormalizzazione. Un dettaglio di quella derivazione si è rivelato decisivo a posteriori: poiché i parametri della GCN non dipendono dalla base di autovettori $U$ (\cref{sec:gcn-kipf-welling}), il modello resta applicabile a un grafo che cambia a ogni finestra, cosa che l'approccio spettrale della prima generazione non consentirebbe (\cref{sec:dynamic-graphs-temporal-question}).

Sul piano empirico la risposta è più cauta, ed è la rassegna del \Cref{ch:literature-review} a imporre la cautela. Il vantaggio rispetto a modelli univariati esiste ed è documentato su più mercati, ma è di entità moderata; sul cripto in particolare l'evidenza si riduce a pochi lavori, non confrontabili tra loro perché condotti su asset, periodi e protocolli diversi (\cref{sec:applications-cryptocurrencies}). La lettura trasversale della \cref{sec:gnn-financial-markets} suggerisce inoltre che la variabilità dei risultati dipenda più dalla scelta del grafo e dal protocollo di valutazione che dalla sofisticazione dell'architettura — il che sposta la domanda interessante da \enquote{quale rete} a \enquote{quale grafo, e valutato come}.

\subsection*{Seconda questione: l'evoluzione della struttura}

Qui l'evidenza è più solida, per la ragione elementare che si tratta di misurare una struttura osservata anziché prevedere qualcosa. La letteratura converge nel documentare che in crisi la correlazione media aumenta, il mercato si compatta e la varianza si concentra su un unico modo dominante (\cref{sec:network-structure-crisis-regimes}); sul cripto il fenomeno si presenta con intensità maggiore, partendo da un livello di correlazione già elevato in condizioni ordinarie.

Il contributo del \Cref{ch:mathematical-foundations} è di rendere questa osservazione misurabile invece che descrittiva. La derivazione della \cref{sec:spectral-analysis}, che identifica ogni autovalore con l'irregolarità del proprio autovettore, dà un significato preciso alla connettività algebrica $\lambda_2$: cresce quando il grafo diventa difficile da separare in gruppi distinti, cioè esattamente quando il mercato smette di articolarsi e si muove in blocco. È una quantità calcolabile da qualunque matrice di correlazione, che non richiede la scelta di una soglia e che tiene conto di come i pesi sono distribuiti sulla struttura, non solo del loro livello medio.

\subsection*{Il punto in cui le due questioni si incontrano}

L'osservazione più utile prodotta da questo lavoro non riguarda però nessuna delle due questioni separatamente, ma la loro relazione: non sono indipendenti, e tirano in direzioni opposte.

Il regime che rende interessante la seconda questione è quello di crisi, in cui il grafo si compatta. È lo stesso regime in cui la prima diventa più difficile: un grafo denso annulla il vantaggio computazionale della convoluzione su grafo e anticipa l'over-smoothing (\cref{sec:known-limitations-gnn}), mentre una soglia di filtraggio calibrata su un periodo tranquillo lascia passare quasi ogni arco proprio quando servirebbe selezionare (\cref{sec:correlation-matrix-to-graph}). Chi progetta un sistema di previsione basato su grafi di correlazione non affronta quindi due problemi separati, ma un compromesso: la struttura è più informativa proprio quando è più difficile da sfruttare. Riconoscerlo in fase di progetto — scegliendo il criterio di filtraggio in funzione del regime, o rinunciando alla profondità in favore di pesi appresi — sembra più promettente che cercare guadagni nell'architettura.

\section{Limiti del lavoro}
\label{sec:limitations}

I limiti che seguono sono già stati riconosciuti nei capitoli in cui sono emersi; raccoglierli qui serve a non lasciarli sparsi.

Il primo è la natura stessa della tesi. Il lavoro è compilativo: non contiene alcuna misura originale, e le affermazioni sulla prima questione di ricerca valgono quanto valgono i lavori recensiti. La \Cref{sec:evaluation-metrics-methodological-pitfalls} ha mostrato che quella letteratura è metodologicamente disomogenea — costi di transazione non dichiarati, numero di configurazioni provate quasi mai riportato, schemi di validazione talvolta inadatti alle serie temporali — per cui la cautela con cui si è formulata la risposta non è una clausola di stile.

Un secondo limite è concettuale e riguarda l'oggetto stesso su cui tutto l'impianto poggia. La correlazione di Pearson misura una relazione lineare, non la dipendenza in generale (\cref{sec:correlation-dependence-pitfalls}), ed è per costruzione una misura di associazione: un arco del grafo può riflettere un fattore comune non osservato — l'esposizione di due token allo stesso settore, o semplicemente al ciclo generale del mercato — invece di una relazione diretta tra i due asset. Vale qui la stessa avvertenza formulata per la causalità di Granger nella \cref{sec:var-baseline}: la struttura appresa o imposta è una regolarità statistica, non un meccanismo.

Terzo, il grafo è imposto e non appreso. Sia i pesi (la correlazione trasformata) sia l'insieme degli archi (il criterio di filtraggio) sono scelte del progettista, e il modello non ha modo di correggerle se sono sbagliate. Il \Cref{ch:from-cnn-to-gnn} ha indicato nel GAT l'alternativa naturale (\cref{sec:message-passing-paradigm}), ma resta un'alternativa non esplorata.

Quarto, e collegato: la trasformazione adottata nella \cref{sec:correlation-matrix-to-graph} ha un costo informativo preciso. La distanza di Mantegna risolve il problema dei pesi negativi e restituisce validità all'apparato spettrale del \Cref{ch:mathematical-foundations}, ma mappa gli asset fortemente anticorrelati in nodi lontani, cioè quasi disconnessi. Una relazione di copertura — quando uno sale, l'altro scende — è informativa quanto una co-variazione, e in questo impianto non viene modellata. Nel mercato cripto la perdita è attenuata dal fatto che le anticorrelazioni forti sono rare, ma è una circostanza empirica favorevole, non una giustificazione della scelta.

Quinto, la non stazionarietà è gestita ma non risolta. La finestra mobile permette di seguire il cambiamento di regime, al prezzo di due parametri — l'ampiezza $T_w$ e la frequenza di aggiornamento del grafo — per i quali la \cref{sec:dynamic-graphs-temporal-question} non ha saputo indicare un criterio soddisfacente, se non il vincolo di rumore derivato dalla legge di Marchenko--Pastur.

Sesto, l'impianto ignora completamente gli attriti di mercato. Costi di transazione, slippage e capienza del mercato non compaiono in alcuna parte della trattazione, e un vantaggio predittivo misurato in RMSE o in accuratezza direzionale può non sopravvivere a costi realistici — tanto più su un mercato che richiede riallineamenti frequenti del portafoglio.

Infine, il perimetro degli asset è una scelta non neutra di cui la tesi non discute le conseguenze. Quali e quante criptovalute costituiscano l'insieme dei nodi $V$ incide sia sui risultati previsivi sia sulle metriche topologiche, e i token meno liquidi hanno serie rumorose che generano archi spuri (\cref{sec:financial-time-series-returns}).

\section{Direzioni future}
\label{sec:future-directions}

Le direzioni più promettenti riguardano, coerentemente con quanto emerso, il grafo più che l'architettura. Si elencano in ordine di distanza crescente dall'impianto discusso qui.

La prima è sostituire i pesi imposti con pesi appresi, cioè adottare un meccanismo di attenzione nel senso del GAT \cite{velickovic2018gat}. La correlazione verrebbe usata solo per decidere quali coppie di asset possano scambiarsi informazione, lasciando al modello il compito di stimare quanto ciascun arco pesi ai fini della previsione. È l'obiezione più naturale all'impianto, e ha il pregio di essere verificabile con una modifica minima.

Un passo oltre c'è l'apprendimento della struttura stessa, non solo dei pesi. Il modello di \emph{neural relational inference} di \textcite{kipf2018nri} inferisce le interazioni tra entità osservando esclusivamente le loro traiettorie, trattando il grafo come variabile latente da stimare insieme alla dinamica. È esattamente il problema della \cref{sec:correlation-matrix-to-graph} riformulato come problema di apprendimento anziché come scelta di progetto, e sarebbe la risposta più radicale al terzo limite della sezione precedente.

Una terza direzione, indipendente dalle prime due, riguarda le feature dei nodi. Nell'impianto discusso ogni nodo porta solo grandezze derivate dalla propria serie di rendimenti, mentre volumi scambiati, indicatori di sentiment estratti da fonti testuali e metriche calcolate sulla blockchain sono informazione disponibile e interamente inutilizzata. È l'estensione meno elegante ma probabilmente la più redditizia, dato che agisce sulla qualità del segnale in ingresso anziché sulla capacità del modello.

Più lontano si colloca l'uso di rappresentazioni alternative del mercato. I grafi di transazione on-chain (\cref{sec:applications-cryptocurrencies}) hanno il vantaggio decisivo di essere osservati e non stimati, e quindi di non soffrire della circolarità tra struttura e dati da prevedere; vivono però a un livello di aggregazione diverso — indirizzi anziché asset — per cui l'ipotesi interessante non è sostituirli al grafo di correlazione, ma combinarli, ad esempio derivandone feature aggregate da associare ai nodi-asset.

Resta, naturalmente, la validazione empirica. Un lavoro sperimentale che confrontasse una GCN con un VAR e con una baseline univariata sugli stessi asset e nello stesso periodo, con protocollo a finestra mobile e metriche multiple riportate insieme alle ipotesi di costo, colmerebbe la lacuna più concreta segnalata dalla \cref{sec:evaluation-metrics-methodological-pitfalls}. Il valore di un lavoro del genere starebbe meno nel risultato che nel protocollo: come la rassegna ha mostrato, in questo campo ciò che scarseggia non sono i modelli, ma i confronti fatti in condizioni comparabili.
